% ============================================================================
% gtopt: High-Performance Generation and Transmission Expansion Planning
%
% Target journal: IEEE Transactions on Power Systems
% Template: IEEEtran (journal mode)
%
% Overleaf-compatible: upload this entire directory as a project.
% To compile locally: latexmk -pdf main.tex
% ============================================================================

\documentclass[journal]{IEEEtran}

% ---- Encoding ----
\usepackage[utf8]{inputenc}
\usepackage[T1]{fontenc}

% ---- Math ----
\usepackage{amsmath,amssymb,amsfonts}

% ---- Graphics and tables ----
\usepackage{graphicx}
\usepackage{booktabs}
\usepackage{multirow}
\usepackage{tabularx}
\usepackage{subcaption}

% ---- Typography ----
\usepackage{textcomp}
\usepackage{xcolor}
\usepackage[binary-units]{siunitx}

% ---- Code listings ----
\usepackage{listings}
\lstset{
  basicstyle=\ttfamily\footnotesize,
  breaklines=true,
  frame=single,
  captionpos=b,
  numbers=left,
  numberstyle=\tiny\color{gray},
  xleftmargin=2em,
  framexleftmargin=1.5em,
}

% ---- Hyperlinks (load last among standard packages) ----
\usepackage[hidelinks,colorlinks=true,citecolor=blue,urlcolor=blue]{hyperref}
\usepackage{cite}

% ---- Algorithms ----
\usepackage{algorithmic}
\usepackage{algorithm}

% ---- Custom commands ----
\newcommand{\gtopt}{\texttt{gtopt}}
\newcommand{\fesop}{FESOP}
\newcommand{\R}{\mathbb{R}}
\newcommand{\Z}{\mathbb{Z}}
\newcommand{\setS}{\mathcal{S}}
\newcommand{\setT}{\mathcal{T}}
\newcommand{\setB}{\mathcal{B}}
\newcommand{\setN}{\mathcal{N}}
\newcommand{\setG}{\mathcal{G}}
\newcommand{\setD}{\mathcal{D}}
\newcommand{\setL}{\mathcal{L}}
\newcommand{\setE}{\mathcal{E}}
\newcommand{\setJ}{\mathcal{J}}
\newcommand{\setR}{\mathcal{R}}
\newcommand{\setW}{\mathcal{W}}
\newcommand{\setH}{\mathcal{H}}
\newcommand{\setZ}{\mathcal{Z}}

% ---- Graphics path ----
\graphicspath{{figures/}{diagrams/}}

% ============================================================================
\begin{document}

\title{%
  \gtopt{}: A High-Performance Open-Source Tool for\\
  Generation and Transmission Expansion Planning%
}

\author{%
  Marcelo~Matus~A.,~\IEEEmembership{Member,~IEEE,}
  Mar\'ia~del~Pilar~Buitrago~Villada%
  % Add more authors as needed:
  % , Author~Three,~\IEEEmembership{Member,~IEEE}
  \thanks{Manuscript received MONTH DAY, 2026; revised MONTH DAY, 2026.}%
  \thanks{M.~Matus~A. is with the Centro Nacional de Inteligencia
    Artificial (CENIA) and the Universidad de Santiago de Chile,
    Santiago, Chile (e-mail: marcelo.matus@usach.cl).}%
  \thanks{M.~P.~Buitrago~Villada is with the Universidad Nacional
    de Colombia, Medell\'in, Colombia.}%
  % \thanks{Digital Object Identifier: 10.1109/TPWRS.20XX.XXXXXXX}%
}

% --- Create title and author block ---
\markboth{IEEE Transactions on Power Systems}%
  {Matus~A. \MakeLowercase{\textit{et al.}}: \gtopt{} for GTEP}

\maketitle

% ========== Abstract ==========
\begin{abstract}
This paper presents \gtopt{}, an open-source, high-performance tool for
Generation and Transmission Expansion Planning (GTEP).  Implemented in modern
C++26, \gtopt{} formulates the joint investment and operational planning
problem as a sparse linear program (LP/MIP) over a multi-stage,
multi-scenario time hierarchy.  The tool supports DC optimal power flow with
Kirchhoff voltage-law constraints, battery energy storage with
charge/discharge efficiency modeling, hydro cascade systems with reservoirs
and turbines, spinning reserve requirements, and capacity expansion with
annualized investment costs.  Three solution methods are provided: a
monolithic LP solver, Stochastic Dual Dynamic Programming (SDDP) with
multi-scenario apertures, and a novel multi-level cascade decomposition.
The architecture features pluggable LP solver backends (CLP, CBC, CPLEX,
HiGHS), native Apache Parquet and CSV I/O for large-scale time-series data,
and a comprehensive Python toolchain including an Excel-to-JSON converter,
PLP and pandapower importers, and web-based graphical interfaces.
Validation against IEEE benchmark networks (4-, 9-, and 14-bus systems)
and cross-comparison with pandapower DC OPF demonstrate correctness, while
computational experiments show significant performance advantages of the
C++ sparse-matrix assembly over Python-based alternatives for systems with
thousands of variables.
\end{abstract}

% ========== Keywords ==========
\begin{IEEEkeywords}
Generation expansion planning, transmission expansion planning, linear
programming, DC optimal power flow, stochastic dual dynamic programming,
battery energy storage, hydro cascade, open-source optimization.
\end{IEEEkeywords}

% ========== Sections ==========
% ============================================================================
% Section I: Introduction
% ============================================================================

\section{Introduction}
\label{sec:introduction}

The global energy transition toward low-carbon power systems demands
sophisticated planning tools that can co-optimize generation, transmission,
and storage investments while capturing the operational complexity of
variable renewable energy sources and hydro cascades.  Generation and
Transmission Expansion Planning (GTEP) is the class of optimization
problems that determines the minimum-cost combination of infrastructure
investments (CAPEX) and dispatch decisions (OPEX) over long planning horizons
under uncertainty~\cite{hemmati2013comprehensive, mahdavi2019tep_survey}.

Traditional GTEP tools---both commercial (PLEXOS, PLP/PROMAX) and
academic---have been implemented in Fortran or commercial modeling
languages, limiting accessibility and extensibility.  The recent emergence
of open-source alternatives such as PyPSA~\cite{brown2018pypsa},
GenX~\cite{jenkins2017genx}, and IDAES-GTEP~\cite{idaes_gtep2024} has
democratized access to planning tools, primarily through Python and Julia
implementations.  However, these interpreted-language tools can face
computational bottlenecks when assembling large sparse LP/MIP problems with
millions of variables and constraints~\cite{frank2016or_opf}.

This paper presents \gtopt{}, a high-performance, open-source GTEP tool
implemented in modern C++26 that bridges the gap between the computational
efficiency of compiled-language solvers and the accessibility of
open-source, JSON/Parquet-based data interfaces.  \gtopt{} descends from
the FESOP (Fabulous Energy System Optimizer) framework~\cite{buitrago2022fesop},
which extended long-term hydrothermal scheduling tools~\cite{pereira2016hydrothermal}
to include capacity expansion and renewable integration.

The main contributions of this work are:

\begin{enumerate}
  \item A complete mathematical formulation of the GTEP problem as a sparse
    LP/MIP, including DC optimal power flow, battery storage, hydro cascades,
    spinning reserves, and multi-stage capacity expansion
    (Section~\ref{sec:formulation}).

  \item Three solution methods---monolithic LP, SDDP with multi-scenario
    apertures, and a novel multi-level cascade decomposition---with formal
    equivalence guarantees (Section~\ref{sec:methods}).

  \item A modular software architecture with pluggable solver backends,
    native Apache Parquet I/O, and companion tools for data preparation
    (Section~\ref{sec:architecture}).

  \item Validation against IEEE benchmark networks (4-bus through 57-bus)
    including battery storage and capacity expansion variants, with
    cross-validation against pandapower DC OPF results
    (Section~\ref{sec:validation}).
\end{enumerate}

The remainder of this paper is organized as follows.
Section~\ref{sec:literature} reviews the state of the art in GTEP tools
and methods.  Section~\ref{sec:formulation} presents the mathematical
formulation.  Section~\ref{sec:methods} describes the three solution
methods.  Section~\ref{sec:architecture} details the software architecture.
Section~\ref{sec:validation} presents validation results on IEEE benchmark
networks.  Section~\ref{sec:conclusions} concludes with a summary and
future research directions.

% ============================================================================
% Section II: Literature Review and State of the Art
% ============================================================================

\section{Literature Review}
\label{sec:literature}

\subsection{Generation and Transmission Expansion Planning}

GTEP is a classical problem in power systems that has been studied
extensively since the 1970s.  Romero et al.~\cite{romero2002test} provided
canonical test systems and mathematical models---including transportation,
hybrid, DC power flow, and disjunctive formulations---that remain standard
benchmarks.  Hemmati et al.~\cite{hemmati2013comprehensive} presented a
comprehensive review covering generation expansion planning (GEP),
transmission expansion planning (TEP), and their integration.  More
recently, Mahdavi et al.~\cite{mahdavi2019tep_survey} classified TEP
literature into categories by mathematical model, solution method, and
uncertainty handling, identifying a trend toward stochastic,
multi-objective, and market-based formulations.

Lumbreras and Ramos~\cite{lumbreras2016progressive} surveyed modern
challenges in TEP including renewable integration, energy storage, and
the need for multi-stage stochastic models.  These challenges motivate
the development of tools that can handle joint generation, transmission,
and storage expansion under operational uncertainty.

\subsection{DC Optimal Power Flow}

The DC power flow approximation linearizes the AC power flow equations
by assuming lossless lines, flat voltage profiles, and small angle
differences~\cite{stott2009dc_opf}.  Despite its simplifications, DC OPF
is the standard formulation for TEP because it yields a linear program
amenable to efficient solution~\cite{frank2016or_opf}.  Molzahn and
Hiskens~\cite{molzahn2019opf_survey} provided a comprehensive survey of
relaxations and approximations, confirming that DC OPF remains the
workhorse for expansion planning despite advances in convex relaxations
of AC OPF.

\gtopt{} adopts the $B$-$\theta$ formulation (susceptance-angle) for
DC OPF, as recommended by Stott et al.~\cite{stott2009dc_opf}, with
extensions for resistive losses via piecewise-linear approximations and
phase-shifting transformers.

\subsection{Stochastic Dual Dynamic Programming}

The Stochastic Dual Dynamic Programming (SDDP) algorithm, introduced by
Pereira and Pinto~\cite{pereira1991sddp}, is the standard method for
solving large-scale multi-stage stochastic programs in hydrothermal
scheduling.  The algorithm constructs piecewise-linear approximations of
the expected cost-to-go function through iterative forward and backward
passes.

Shapiro~\cite{shapiro2011sddp_analysis} provided theoretical convergence
analysis of SDDP, and Philpott and de~Matos~\cite{philpott2013sddp_convergence}
extended the algorithm with risk-averse sampling.  Recent advances include
batch learning SDDP~\cite{leclercq2023batch_sddp}, which frames the
algorithm as Q-learning for improved convergence, and extensions to
stochastic OPF under uncertainty~\cite{fletcher2024sddp_opf}.

In the Julia ecosystem, Dowson and Kapelevich~\cite{dowson2020sddp_jl}
developed SDDP.jl, a modular SDDP implementation.  \gtopt{} provides a
C++ SDDP implementation with support for multi-scenario apertures, elastic
filters for feasibility management, and boundary cost cuts for rolling
horizon studies.

\subsection{Energy Storage in Expansion Planning}

The integration of battery energy storage systems (BESS) into expansion
planning has received significant attention.  Go et al.~\cite{Go2016ess_expansion}
assessed the economic value of co-optimized grid-scale storage investments.
Frazier et al.~\cite{frazier2023bess_milp} presented an integrated MILP
for generation, transmission, and storage expansion with Benders
decomposition.

Sandia National Laboratories released QuESt Planning~\cite{sandia2024quest},
an open-source capacity expansion tool focused on storage, while
IDAES-GTEP~\cite{idaes_gtep2024} provides a modular Pyomo-based framework
for GTEP.  \gtopt{} supports BESS with detailed charge/discharge efficiency
modeling, state-of-charge tracking across stages, and co-located
generation-storage coupling via its converter mechanism.

\subsection{Open-Source Power System Tools}

Table~\ref{tab:tool_comparison} compares \gtopt{} with the leading
open-source GTEP tools.  PyPSA~\cite{brown2018pypsa} is the most widely
adopted, offering multi-sector coupling and a large community, but uses
Python for LP assembly which can bottleneck on large problems.
GenX~\cite{jenkins2017genx} provides a flexible Julia/JuMP framework but
lacks native DC OPF and hydro cascade modeling.  pandapower~\cite{thurner2018pandapower}
excels at power flow analysis but does not support expansion planning.

\begin{table*}[t]
  \centering
  \caption{Comparison of Open-Source GTEP and Power System Tools}
  \label{tab:tool_comparison}
  \begin{tabular}{@{}lcccccc@{}}
    \toprule
    \textbf{Feature} & \textbf{gtopt} & \textbf{PyPSA} & \textbf{GenX}
      & \textbf{IDAES-GTEP} & \textbf{QuESt} & \textbf{PLP} \\
    \midrule
    Language          & C++26  & Python  & Julia   & Python  & Python  & Fortran \\
    License           & BSD-3  & MIT     & GPL-2   & BSD-3   & BSD-3   & Proprietary \\
    DC OPF (Kirchhoff)& \checkmark & \checkmark & Transport & \checkmark & -- & \checkmark \\
    Hydro cascades    & \checkmark & Partial & -- & -- & -- & \checkmark \\
    Battery storage   & \checkmark & \checkmark & \checkmark & Partial & \checkmark & Partial \\
    SDDP solver       & \checkmark & -- & -- & -- & -- & \checkmark \\
    Capacity expansion& \checkmark & \checkmark & \checkmark & \checkmark & \checkmark & -- \\
    Multi-scenario    & \checkmark & Partial & Partial & -- & -- & \checkmark \\
    Multi-sector      & -- & \checkmark & -- & -- & -- & -- \\
    Parquet I/O       & \checkmark & -- & -- & -- & -- & -- \\
    Web GUI           & \checkmark & -- & -- & -- & \checkmark & -- \\
    \bottomrule
  \end{tabular}
\end{table*}

\gtopt{} distinguishes itself through: (i)~C++ sparse-matrix assembly
for high performance, (ii)~complete hydro cascade modeling with
reservoirs, waterways, turbines, and filtration, (iii)~three solution
methods including SDDP and cascade decomposition, (iv)~native Apache
Parquet I/O for large time-series, and (v)~a comprehensive Python
toolchain for data preparation and cross-validation.

% ============================================================================
% Section III: Mathematical Formulation
% ============================================================================

\section{Mathematical Formulation}
\label{sec:formulation}

\subsection{Sets and Indices}

The GTEP problem is defined over the following sets:

\begin{itemize}
  \item $\setS$: scenarios, indexed by $s$, with probability $\pi_s$
  \item $\setT$: stages (planning periods), indexed by $t$
  \item $\setB_t$: blocks (operating periods) in stage $t$, indexed by $b$,
        with duration $\Delta_b$ [h]
  \item $\setN$: buses (electrical nodes), indexed by $n$
  \item $\setG$: generators, indexed by $g$, with bus $n_g$
  \item $\setD$: demands, indexed by $d$, with bus $n_d$
  \item $\setL$: transmission lines, indexed by $\ell$, connecting buses
        $a_\ell$ and $b_\ell$
  \item $\setE$: batteries (energy storage), indexed by $e$
  \item $\setJ$: hydraulic junctions, indexed by $j$
  \item $\setR$: reservoirs, indexed by $r$, with junction $j_r$
  \item $\setW$: waterways, indexed by $w$
  \item $\setH$: turbines, indexed by $h$
  \item $\setZ$: reserve zones, indexed by $z$
\end{itemize}

\begin{figure}[t]
  \centering
  \begin{tikzpicture}[
    scenbox/.style={draw, rounded corners=6pt, fill=blue!5,
      minimum width=7.6cm, inner sep=4pt},
    stagebox/.style={draw, rounded corners=4pt, fill=blue!10,
      minimum width=3.4cm, inner sep=3pt},
    blockbox/.style={draw, rounded corners=2pt, fill=blue!18,
      minimum width=0.55cm, minimum height=0.55cm,
      font=\scriptsize, inner sep=1pt},
    lbl/.style={font=\scriptsize\itshape, text=black!70},
  ]
    % --- Scenario 1 ---
    \node[blockbox] (s1t1b1) at (0,0) {$b_1$};
    \node[blockbox, right=1pt of s1t1b1] (s1t1b2) {$b_2$};
    \node[blockbox, right=1pt of s1t1b2] (s1t1b3) {$b_3$};
    \node[stagebox, fit=(s1t1b1)(s1t1b3), label={[lbl]above:{Stage $t{=}1$, $\delta_1$}}] (s1t1) {};

    \node[blockbox, right=12pt of s1t1b3] (s1t2b1) {$b_1$};
    \node[blockbox, right=1pt of s1t2b1] (s1t2b2) {$b_2$};
    \node[blockbox, right=1pt of s1t2b2] (s1t2b3) {$b_3$};
    \node[blockbox, right=1pt of s1t2b3] (s1t2b4) {$b_4$};
    \node[stagebox, fit=(s1t2b1)(s1t2b4), label={[lbl]above:{Stage $t{=}2$, $\delta_2$}}] (s1t2) {};

    \node[scenbox, fit=(s1t1)(s1t2),
      label={[lbl]above:{Scenario $s{=}1$, $\pi_1$}}] (sc1) {};

    % --- Scenario 2 ---
    \node[blockbox, below=28pt of s1t1b1] (s2t1b1) {$b_1$};
    \node[blockbox, right=1pt of s2t1b1] (s2t1b2) {$b_2$};
    \node[blockbox, right=1pt of s2t1b2] (s2t1b3) {$b_3$};
    \node[stagebox, fit=(s2t1b1)(s2t1b3), label={[lbl]above:{Stage $t{=}1$}}] (s2t1) {};

    \node[blockbox, right=12pt of s2t1b3] (s2t2b1) {$b_1$};
    \node[blockbox, right=1pt of s2t2b1] (s2t2b2) {$b_2$};
    \node[blockbox, right=1pt of s2t2b2] (s2t2b3) {$b_3$};
    \node[blockbox, right=1pt of s2t2b3] (s2t2b4) {$b_4$};
    \node[stagebox, fit=(s2t2b1)(s2t2b4), label={[lbl]above:{Stage $t{=}2$}}] (s2t2) {};

    \node[scenbox, fit=(s2t1)(s2t2),
      label={[lbl]above:{Scenario $s{=}2$, $\pi_2$}}] (sc2) {};

    % --- Duration label ---
    \draw[decorate, decoration={brace, amplitude=3pt, mirror}]
      (s1t1b1.south west) -- (s1t1b1.south east)
      node[midway, below=3pt, font=\tiny] {$\Delta_b$};
  \end{tikzpicture}
  \caption{Time hierarchy: scenarios $\setS$ contain stages $\setT$,
    each with blocks $\setB_t$.  Probabilities $\pi_s$, discount factors
    $\delta_t$, and block durations $\Delta_b$ weight the objective.}
  \label{fig:time_hierarchy}
\end{figure}

\subsection{Decision Variables}

\noindent\textbf{Operational variables} (per scenario $s$, stage $t$, block $b$):
\begin{align}
  p_{g,s,t,b}  &\in [\underline{p}_g, \overline{p}_g]
    & \text{generator output [MW]} \label{eq:pgen} \\
  \ell_{d,s,t,b} &\in [0, \overline{\ell}_d]
    & \text{served demand [MW]} \label{eq:load} \\
  f_{d,s,t,b}  &\geq 0
    & \text{unserved demand [MW]} \label{eq:fail} \\
  \phi_{\ell,s,t,b} &\in [-\overline{\phi}^{ba}_\ell, \overline{\phi}^{ab}_\ell]
    & \text{line flow [MW]} \label{eq:flow} \\
  \theta_{n,s,t,b} &\in \R
    & \text{voltage angle [rad]} \label{eq:theta}
\end{align}

\noindent\textbf{Storage state variables}:
\begin{align}
  e_{e,s,t,b}  &\in [\underline{e}_e, \overline{e}_e]
    & \text{battery SoC [MWh]} \label{eq:soc} \\
  v_{r,s,t,b}  &\in [\underline{v}_r, \overline{v}_r]
    & \text{reservoir volume [hm\textsuperscript{3}]} \label{eq:vol}
\end{align}

\noindent\textbf{Investment variables} (per stage $t$):
\begin{equation}
  m_{g,t} \in [0, \overline{m}_g]
  \quad \text{expansion modules built} \label{eq:expmod}
\end{equation}

\subsection{Objective Function}

Minimize total expected discounted cost:
\begin{equation}
\label{eq:objective}
\min \sum_{s \in \setS} \pi_s \sum_{t \in \setT} \delta_t
\left[
  \sum_{b \in \setB_t} \Delta_b \cdot C^\text{op}_{s,t,b}
  + C^\text{inv}_t
\right]
\end{equation}
where $\delta_t = (1+r)^{-\tau_t/8760}$ is the discount factor with annual
rate $r$ and $\tau_t$ the cumulative hours to stage $t$.

The operational cost for each $(s,t,b)$ is:
\begin{equation}
\label{eq:opcost}
C^\text{op}_{s,t,b} = \sum_{g \in \setG} c_g \, p_{g,s,t,b}
  + \sum_{d \in \setD} c^f_d \, f_{d,s,t,b}
  + \sum_{\ell \in \setL} c^\phi_\ell \, |\phi_{\ell,s,t,b}|
\end{equation}

The investment cost for stage $t$ is:
\begin{equation}
\label{eq:invcost}
C^\text{inv}_t = \sum_{g \in \setG} \kappa_g \, m_{g,t}
  + \sum_{d \in \setD} \kappa_d \, m_{d,t}
  + \sum_{\ell \in \setL} \kappa_\ell \, m_{\ell,t}
  + \sum_{e \in \setE} \kappa_e \, m_{e,t}
\end{equation}
where $c_g$ is the variable generation cost [\$/MWh], $c^f_d$ is the
demand failure penalty [\$/MWh], $c^\phi_\ell$ is the transmission cost
[\$/MWh], and $\kappa_{(\cdot)}$ is the annualized investment cost
[\$/year] for each component.  All objective coefficients are divided by a
scaling factor $\sigma$ (\texttt{scale\_objective}, default 1000) to
improve solver numerical conditioning.  Concretely, each LP cost
coefficient is weighted by $\omega_{s,t,b} = \pi_s \delta_t \Delta_b /
\sigma$, so the reported objective equals $z/\sigma$.

\subsection{Bus Power Balance}

At every bus $n$, for each $(s,t,b)$:
\begin{multline}
\label{eq:balance}
\sum_{\substack{g \in \setG \\ n_g = n}} (1{-}\lambda_g) p_{g,s,t,b}
- \sum_{\substack{d \in \setD \\ n_d = n}}
  (1{+}\lambda_d)(\ell_{d,s,t,b} + f_{d,s,t,b}) \\
+ \sum_{\ell \in \setL} A_{n,\ell} \, \phi_{\ell,s,t,b} = 0
\end{multline}
where $\lambda_g, \lambda_d$ are injection and withdrawal loss factors,
respectively.  The line flow terms use directional contributions rather
than a simple incidence matrix: for each line~$\ell$ from bus~$a_\ell$
to~$b_\ell$, the sending bus sees $-\phi^+_\ell$ (or $+(1-\lambda_\ell)
\phi^-_\ell$ when losses are modeled) and the receiving bus sees
$+(1-\lambda_\ell)\phi^+_\ell$ (or $-\phi^-_\ell$), where $\lambda_\ell$
is the line loss factor.  When $\lambda_\ell=0$, a single bidirectional
variable suffices and the incidence reduces to $A_{n,\ell} \in
\{+1,-1,0\}$.  The dual of~\eqref{eq:balance} is the locational marginal
price (LMP) at bus~$n$.

\subsection{DC Power Flow (Kirchhoff Constraints)}

When \texttt{use\_kirchhoff} is enabled, for each line $\ell$ and
$(s,t,b)$:
\begin{equation}
\label{eq:kirchhoff}
\phi_{\ell,s,t,b} = \frac{V_\ell^2}{X_\ell}
  \left(\theta_{a_\ell,s,t,b} - \theta_{b_\ell,s,t,b}\right)
\end{equation}
where $X_\ell$ is the line reactance [$\Omega$] and $V_\ell$ is the
nominal voltage [kV].  A scaling factor $\sigma_\theta$
(\texttt{scale\_theta}, default 1000) is applied to the angle variables
for numerical conditioning: $\theta^\text{scaled} = \sigma_\theta \cdot
\theta$.

Line flow is bounded by transfer capacities:
\begin{equation}
-\overline{\phi}^{ba}_\ell \leq \phi_{\ell,s,t,b}
  \leq \overline{\phi}^{ab}_\ell
\end{equation}

\paragraph{Transmission losses.}
When line losses are enabled ($\lambda_\ell > 0$), separate forward and
reverse flow variables $\phi^+_\ell, \phi^-_\ell \geq 0$ replace the
single bidirectional variable.  The receiving bus sees only
$(1-\lambda_\ell)$ of the dispatched power, so losses are allocated at
the receiving end.  The Kirchhoff constraint~\eqref{eq:kirchhoff} then
uses $\phi^+_\ell - \phi^-_\ell$ in place of the single~$\phi_\ell$.

\paragraph{Transformer model.}
For tap-changing and phase-shifting transformers, the Kirchhoff
constraint generalizes to:
\begin{equation}
\label{eq:transformer}
-\theta_{a_\ell} + \theta_{b_\ell}
  + (\tau_\ell \chi_\ell)(\phi^+_\ell - \phi^-_\ell)
  = -\sigma_\theta \, \varphi_{\ell,\text{rad}}
\end{equation}
where $\tau_\ell$ is the off-nominal tap ratio (default~1),
$\chi_\ell = \sigma_\theta X_\ell / V_\ell^2$ is the scaled reactance,
and $\varphi_{\ell,\text{rad}}$ is the phase-shift angle in radians.
When $\tau_\ell = 1$ and $\varphi_\ell = 0$,~\eqref{eq:transformer}
reduces to~\eqref{eq:kirchhoff}.

\subsection{Battery Energy Storage}

State-of-charge balance for battery $e$ between consecutive blocks:
\begin{equation}
\label{eq:soc_balance}
e_{e,s,t,b} = e_{e,s,t,b-1}
  + p^\text{in}_{e,s,t,b} \cdot \eta^\text{in}_e \cdot \Delta_b
  - p^\text{out}_{e,s,t,b} \cdot \Delta_b / \eta^\text{out}_e
\end{equation}
where $\eta^\text{in}_e$ and $\eta^\text{out}_e$ are charge and discharge
efficiencies, and $p^\text{in}, p^\text{out}$ are charge and discharge
power [MW].  Initial SoC $e_{e,\cdot,\cdot,0} = e^\text{ini}_e$ is an
equality constraint in the first stage.

\begin{figure}[t]
  \centering
  \begin{tikzpicture}[
    >=Stealth,
    axis/.style={->, thick},
    charge/.style={fill=green!15},
    discharge/.style={fill=red!10},
  ]
    % Axes
    \draw[axis] (0,0) -- (7.2,0)
      node[right, font=\scriptsize] {Block (hour)};
    \draw[axis] (0,0) -- (0,3.2)
      node[above, font=\scriptsize] {SoC (\%)};

    % Y-axis ticks
    \foreach \y/\lbl in {0/0, 0.6/20, 1.2/40, 1.8/60, 2.4/80, 3.0/100} {
      \draw (-0.08,\y) -- (0.08,\y)
        node[left=2pt, font=\tiny] {\lbl};
    }
    % X-axis ticks
    \foreach \x/\lbl in
      {0.3/1, 1.5/6, 2.7/10, 3.6/13, 4.5/16, 5.7/20, 6.9/24} {
      \draw (\x,-0.08) -- (\x,0.08)
        node[below=2pt, font=\tiny] {\lbl};
    }

    % SoC data points
    \coordinate (p0) at (0.0, 0.9);   % hour 0:  30%
    \coordinate (p1) at (0.9, 0.6);   % hour 3:  20%
    \coordinate (p2) at (1.8, 0.6);   % hour 6:  20%
    \coordinate (p3) at (2.7, 1.5);   % hour 10: 50%
    \coordinate (p4) at (3.6, 2.7);   % hour 13: 90%
    \coordinate (p5) at (4.2, 2.7);   % hour 15: 90%
    \coordinate (p6) at (5.1, 1.5);   % hour 18: 50%
    \coordinate (p7) at (6.0, 0.9);   % hour 21: 30%
    \coordinate (p8) at (6.9, 0.9);   % hour 24: 30%

    % Charge shading (hours 6--15)
    \fill[charge] (p2) -- (p3) -- (p4) -- (p5)
      -- (p5 |- 0,0) -- (p2 |- 0,0) -- cycle;

    % Discharge shading (hours 15--21)
    \fill[discharge] (p5) -- (p6) -- (p7)
      -- (p7 |- 0,0) -- (p5 |- 0,0) -- cycle;

    % SoC line
    \draw[thick, blue!80] (p0) -- (p1) -- (p2) -- (p3) -- (p4)
      -- (p5) -- (p6) -- (p7) -- (p8);

    % Labels
    \node[font=\tiny, green!50!black, above] at (3.15,2.1)
      {charge};
    \node[font=\tiny, red!60!black, above] at (4.95,1.8)
      {discharge};

    % Solar peak annotation
    \draw[{Stealth[length=3pt]}-, thin, gray]
      (3.6,2.7) -- (4.3,3.1)
      node[right, font=\tiny] {solar peak};

    % Evening peak annotation
    \draw[{Stealth[length=3pt]}-, thin, gray]
      (5.1,1.5) -- (5.8,2.2)
      node[right, font=\tiny] {evening peak};
  \end{tikzpicture}
  \caption{Battery SoC over 24~blocks: charge during solar hours
    (green), discharge during evening peak (red).}
  \label{fig:battery_soc}
\end{figure}

\subsection{Converter Constraints}

A converter~$v$ couples a battery~$e$ to the electrical network through
a generator~$g$ (discharge path) and a demand~$d$ (charge path) via a
power conversion rate~$\rho_v$:
\begin{align}
\label{eq:conv_discharge}
p_{g,s,t,b} &= \rho_v \cdot p^\text{out}_{e,s,t,b} \\
\label{eq:conv_charge}
\ell_{d,s,t,b} &= \rho_v \cdot p^\text{in}_{e,s,t,b}
\end{align}
When the converter has a capacity limit~$\overline{C}_v$, a joint bound
applies: $p_{g,s,t,b} + \ell_{d,s,t,b} \leq \overline{C}_{v,t}$.
Converter capacity is expandable using the same modular framework
as~\eqref{eq:capacity}.

\subsection{Hydro Cascade}

For each junction $j$ with reservoir $r$, the water balance is:
\begin{equation}
\label{eq:hydro_balance}
v_{r,s,t,b} = v_{r,s,t,b-1}
  + \alpha_r \Delta_b \!\left(
    \textstyle\sum_{w \in \text{in}(j)} q_{w}
    - \textstyle\sum_{w \in \text{out}(j)} q_{w}
    + F_j
  \right)
\end{equation}
where $q_w$ is the waterway discharge [m$^3$/s], $F_j$ is the exogenous
inflow at junction $j$, and $\alpha_r = 0.0036$ converts flow rates to
volume increments: since 1~m$^3$/s flowing for 1~hour equals
3600~m$^3$ = 0.0036~hm$^3$, we have $\alpha_r = 3.6 \times 10^{-3}$
[hm$^3$/(m$^3$/s$\cdot$h)].

Each turbine $h$ linking waterway $w$ to generator $g$ provides:
\begin{equation}
\label{eq:turbine}
p_{g,s,t,b} \leq \rho_h \cdot q_{w,s,t,b}
\end{equation}
where $\rho_h$ is the water-to-power conversion rate [MW$\cdot$s/m$^3$].

\paragraph{Volume-dependent turbine efficiency.}
In practice, the conversion rate depends on the hydraulic head, which
varies with reservoir volume.  When defined, $\rho_h$ is replaced by a
concave piecewise-linear function:
\begin{equation}
\label{eq:vol_efficiency}
\rho_h(v) = \min_{i=1}^{N}
  \bigl\{ c_i + m_i (v - v_i) \bigr\}
\end{equation}
where $(v_i, m_i, c_i)$ are volume breakpoints, slopes, and intercepts
of~$N$ segments.  In the SDDP method, the active segment is updated at
each forward-pass iteration using the current reservoir volume.

\paragraph{Filtration (seepage).}
Seepage from a waterway into a reservoir is modeled as a linear function
of the average volume:
\begin{equation}
\label{eq:filtration}
\varphi_{i,s,t,b} = a_i + b_i \cdot
  \frac{v_{r,s,t,b-1} + v_{r,s,t,b}}{2}
\end{equation}
where $a_i$ is a constant term and $b_i$ is the seepage rate per unit
volume.  As with turbine efficiency, this admits a piecewise-linear
extension selected by the last known volume.

\begin{figure}[t]
  \centering
  \begin{tikzpicture}[
    res/.style={trapezium, trapezium angle=75, draw, fill=cyan!15,
      minimum width=1.4cm, minimum height=0.7cm, font=\scriptsize,
      trapezium stretches body},
    junc/.style={circle, draw, fill=gray!20, minimum size=0.4cm,
      font=\tiny, inner sep=1pt},
    turb/.style={regular polygon, regular polygon sides=3, draw,
      fill=yellow!25, minimum size=0.5cm, font=\tiny, inner sep=0pt},
    gen/.style={circle, draw, fill=orange!20, minimum size=0.5cm,
      font=\tiny},
    flow/.style={-{Stealth[length=4pt]}, thick, blue!70},
    spill/.style={-{Stealth[length=4pt]}, thick, red!50, dashed},
  ]
    % Reservoir 1
    \node[res] (r1) at (0,2.2) {$R_1$};
    \node[font=\tiny\itshape, text=black!70, right=1pt of r1] {$v_1$};

    % Junction 1
    \node[junc] (j1) at (0,1.1) {$j_1$};

    % Inflow
    \draw[flow] (-1.2,2.8) -- (j1.north west)
      node[midway, left, font=\tiny] {inflow $F_1$};

    % R1 -> J1
    \draw[flow] (r1) -- (j1);

    % Turbine
    \node[turb] (t1) at (1.2,0.6) {};
    \node[font=\tiny, right=1pt of t1] {$h_1$};

    % Junction 2
    \node[junc] (j2) at (0,-0.2) {$j_2$};

    % J1 -> T1
    \draw[flow] (j1) -- (t1) node[midway, above, font=\tiny] {$q_w$};

    % T1 -> Generator
    \node[gen] (g1) at (2.5,0.6) {\raisebox{-0.5pt}{\scalebox{0.5}{$\sim$}}};
    \node[font=\tiny, right=1pt of g1] {gen};
    \draw[-{Stealth[length=4pt]}, thick, orange!70] (t1) -- (g1);

    % J1 -> J2 spillway
    \draw[spill] (j1) -- (j2) node[midway, left, font=\tiny] {spill};

    % Reservoir 2
    \node[res] (r2) at (2,-0.2) {$R_2$};
    \node[font=\tiny\itshape, text=black!70, right=1pt of r2] {$v_2$};

    % Junction 3
    \node[junc] (j3) at (1,-1.2) {$j_3$};

    % J2 -> R2 waterway
    \draw[flow] (j2) -- (r2);

    % T1 -> J3 (tailwater)
    \draw[flow] (t1) |- (j3);

    % R2 -> J3
    \draw[flow] (r2) -- (j3);

    % Outflow
    \draw[flow] (j3) -- (1,-2.0) node[below, font=\tiny] {outflow};
  \end{tikzpicture}
  \caption{Hydro cascade example: two reservoirs ($R_1$, $R_2$),
    three junctions ($j_1$--$j_3$), one turbine ($h_1$) driving a
    generator, with inflow, spillway, and waterway connections.}
  \label{fig:hydro_cascade}
\end{figure}

\subsection{Capacity Expansion}

Installed capacity evolves as:
\begin{equation}
\label{eq:capacity}
C_{g,t} = C_{g,t-1}(1-\xi_g) + M_g \cdot m_{g,t}
\end{equation}
where $C_{g,0}$ is the initial installed capacity, $\xi_g$ is the annual
derating factor (\texttt{annual\_derating}), $M_g$ is the capacity per
expansion module (\texttt{expcap}), and $m_{g,t}$ is the number of modules
built.  The generator output is bounded by the installed capacity:
\begin{equation}
p_{g,s,t,b} \leq C_{g,t}
\end{equation}
Analogous expansion constraints apply to demands, lines, batteries, and
converters.

\subsection{Generator Profile Constraints}

When a generator has a capacity profile $\phi_g(s,t,b)$ (e.g., solar
irradiance or wind availability), the output is forced to follow the
profile:
\begin{equation}
\label{eq:profile}
p_{g,s,t,b} + s_{g,s,t,b} = \phi_g(s,t,b) \cdot C_{g,t}
\end{equation}
where $\phi_g(s,t,b) \in [0,1]$ is the per-block capacity factor and
$s_{g,s,t,b} \geq 0$ is the curtailed (spilled) generation.  The
spillover variable carries cost $c_g^\text{spill} \cdot \omega_{s,t,b}$,
which may be set to zero to allow free curtailment of renewables.  This
lets the optimizer spill excess renewable production when it is not
economically efficient to deliver.

\subsection{Demand Minimum Energy}

An optional minimum energy constraint ensures that a demand~$d$ is
served a minimum total energy within each stage:
\begin{equation}
\label{eq:min_energy}
\sum_{b \in \setB_t} \ell^\text{man}_{d,s,t,b} \cdot \Delta_b
  \geq E_d^\text{min}
\end{equation}
where $\ell^\text{man}_{d,s,t,b}$ is the mandatory served load and
$E_d^\text{min}$ [MWh] is the stage energy requirement.

\subsection{Reserve Constraints}

For each reserve zone $z$ and $(s,t,b)$:
\begin{equation}
\label{eq:reserve_up}
\sum_{g \in \mathcal{P}_z} u^\text{up}_{g,s,t,b}
  + u^\text{fail,up}_{z,s,t,b}
  \geq R^\text{up}_{z,s,t,b}
\end{equation}
where $u^\text{up}_{g,s,t,b}$ is the up-reserve provision from generator
$g$, $u^\text{fail,up}_{z,s,t,b}$ is unserved reserve, and
$R^\text{up}_{z,s,t,b}$ is the zone requirement.  A symmetric constraint
exists for down-reserve.  Reserve provision is bounded by the generator's
headroom: $u^\text{up}_{g,s,t,b} \leq C_{g,t} - p_{g,s,t,b}$.

\subsection{Constraint Summary}

Table~\ref{tab:constraints} lists all constraint families in the
formulation.

\begin{table}[t]
\centering
\caption{Complete constraint summary.}
\label{tab:constraints}
\small
\begin{tabular}{@{}clll@{}}
\toprule
\textbf{\#} & \textbf{Constraint} & \textbf{Eq.} & $\boldsymbol{\forall}$ \\
\midrule
C1  & Bus power balance             & \eqref{eq:balance}       & $n,s,t,b$ \\
C2  & Kirchhoff voltage law         & \eqref{eq:kirchhoff}     & $\ell,s,t,b$ \\
C3  & Generator output bounds       & \eqref{eq:pgen}          & $g,s,t,b$ \\
C4  & Capacity expansion            & \eqref{eq:capacity}      & $g,t$ \\
C5  & Generator profile             & \eqref{eq:profile}       & $g,s,t,b$ \\
C6  & Demand min.\ energy           & \eqref{eq:min_energy}    & $d,s,t$ \\
C7  & Line capacity                 & ---                      & $\ell,s,t,b$ \\
C8  & Battery SoC                   & \eqref{eq:soc_balance}   & $e,s,t,b$ \\
C9  & Converter coupling            & \eqref{eq:conv_discharge}& $v,s,t,b$ \\
C10 & Reserve requirements          & \eqref{eq:reserve_up}    & $z,s,t,b$ \\
C11 & Reserve--generator coupling   & ---                      & $p,s,t,b$ \\
C12 & Hydro water balance           & \eqref{eq:hydro_balance} & $j,s,t,b$ \\
C13 & Turbine conversion            & \eqref{eq:turbine}       & $h,s,t,b$ \\
C14 & Transformer / PST             & \eqref{eq:transformer}   & $\ell,s,t,b$ \\
\bottomrule
\end{tabular}
\end{table}

% ============================================================================
% Section IV: Solution Methods
% ============================================================================

\section{Solution Methods}
\label{sec:methods}

\gtopt{} implements three complementary solution methods, selectable via the
\texttt{method} option.  All three solve the same underlying GTEP formulation
(Section~\ref{sec:formulation}); they differ in how they decompose the
problem across the time and scenario dimensions.

\subsection{Monolithic LP}
\label{sec:monolithic}

The default method assembles the entire multi-scenario, multi-stage problem
into a single sparse LP/MIP and passes it to the solver backend.  The LP
is constructed in compressed sparse column (CSC) format via the
\texttt{FlatLinearProblem} class.

For a system with $N$ buses, $G$ generators, $D$ demands, $L$ lines,
$E$ batteries, $B$ blocks per stage, $T$ stages, and $S$ scenarios, the
approximate problem size is:

\begin{table}[h]
  \centering
  \caption{LP Problem Size Estimates}
  \label{tab:lp_size}
  \begin{tabular}{@{}ll@{}}
    \toprule
    Quantity & Approximate count \\
    \midrule
    Operational variables & $(G{+}2D{+}2L{+}3E{+}N) \times S \times T \times B$ \\
    Investment variables  & $(G{+}D{+}L{+}E) \times T$ \\
    Bus balance rows      & $N \times S \times T \times B$ \\
    Kirchhoff rows        & $L \times S \times T \times B$ (if DC OPF) \\
    SoC balance rows      & $E \times S \times T \times B$ \\
    Capacity rows         & $(G{+}D{+}L{+}E) \times T$ \\
    \bottomrule
  \end{tabular}
\end{table}

The monolithic method is exact (no decomposition error) and is recommended
for problems where the total number of LP variables fits in memory.  For
the IEEE~57-bus system with 24 hourly blocks, the LP has approximately
$57 \times 24 \times (7 + 84 + 160 + 57) \approx 420\,000$ variables.

\subsection{Stochastic Dual Dynamic Programming (SDDP)}
\label{sec:sddp}

For multi-stage stochastic problems with many scenarios, \gtopt{} implements
SDDP following the algorithm of Pereira and
Pinto~\cite{pereira1991sddp}.  The stages are decomposed into
individual subproblems connected by Benders optimality cuts that
approximate the expected future cost.

\begin{algorithm}[h]
  \caption{SDDP Algorithm in \gtopt{}}
  \label{alg:sddp}
  \begin{algorithmic}[1]
    \STATE Initialize: empty cut sets $\mathcal{C}_t = \emptyset$
      for all $t$
    \FOR{iteration $k = 1, 2, \ldots, K$}
      \STATE \textbf{Forward pass}: Solve stages $t=1,\ldots,T$
        sequentially for sampled scenarios; record state variables
        (SoC, reservoir volumes)
      \STATE \textbf{Backward pass}: For $t=T,\ldots,1$, solve
        each stage for all aperture scenarios; compute dual
        variables; generate Benders cuts:
        \[
          \alpha_t \geq \hat{\alpha}_t + \hat{\pi}_t^\top
            (x_t - \hat{x}_t)
        \]
      \STATE Add cuts to $\mathcal{C}_t$
      \STATE Check convergence: upper/lower bound gap $\leq \epsilon$
    \ENDFOR
  \end{algorithmic}
\end{algorithm}

\gtopt{}'s SDDP implementation includes several enhancements:

\begin{itemize}
  \item \textbf{Apertures}: configurable backward-pass scenario sampling
    that groups scenarios into subsets for cut generation, reducing
    computational cost while preserving solution quality.
  \item \textbf{Elastic filters}: infeasible subproblems are handled by
    relaxing state variable bounds with penalty costs, ensuring the
    algorithm does not stall.
  \item \textbf{Boundary cuts}: pre-computed future cost approximations
    at the terminal stage, enabling rolling horizon studies.
  \item \textbf{Hot start}: cuts from previous runs can be loaded to
    accelerate convergence.
  \item \textbf{Cut management}: dominated cuts are periodically purged
    to control LP size growth.
\end{itemize}

\subsection{Cascade Decomposition}
\label{sec:cascade}

The cascade method is a novel multi-level hybrid that decomposes the
problem hierarchically:

\begin{enumerate}
  \item \textbf{Level~0 (strategic)}: A coarsened version of the problem
    (aggregated blocks, reduced scenarios) is solved monolithically to
    determine investment decisions and approximate cost-to-go functions.
  \item \textbf{Level~1 (tactical)}: SDDP is applied to the full
    stochastic problem using the Level~0 investments as warm-start and
    the cost-to-go cuts as initial approximations.
  \item \textbf{Level~2+ (operational)}: Additional refinement levels
    can solve detailed operational subproblems with fixed investments.
\end{enumerate}

The cascade method inherits SDDP's convergence guarantees while using the
monolithic coarse solution to dramatically reduce the number of SDDP
iterations required.  It is particularly effective for long-horizon problems
(10+ stages) with many scenarios.

\subsection{Solver Backend Architecture}

All three methods delegate LP/MIP solution to a pluggable solver backend
loaded as a shared library at runtime.  Currently supported backends are:

\begin{itemize}
  \item \textbf{CLP/CBC} (COIN-OR): Default open-source backend.  CLP for
    LP, CBC for MIP~\cite{forrest2005cbc}.
  \item \textbf{CPLEX} (IBM): Commercial high-performance solver, accessed
    via the C Callable Library.
  \item \textbf{HiGHS}: Open-source solver with parallel simplex and
    interior-point methods~\cite{huangfu2018highs}.
\end{itemize}

The backend is auto-detected by priority (CPLEX $>$ HiGHS $>$ CBC $>$ CLP)
or explicitly selected via the \texttt{--solver} flag.  Solver-specific
options (algorithm, threads, tolerances, time limits) are configured
through the \texttt{solver\_options} JSON section.

% ============================================================================
% Section V: Software Architecture
% ============================================================================

\section{Software Architecture}
\label{sec:architecture}

\subsection{Design Principles}

\gtopt{} is designed around four core principles:

\begin{enumerate}
  \item \textbf{Performance}: C++26 with flat-map-based sparse matrices
    for fast LP assembly.  The \texttt{boost::container::flat\_map}
    (aliased as \texttt{gtopt::flat\_map}) provides 10--27$\times$ faster
    iteration than \texttt{std::map} in LP coefficient assembly.
  \item \textbf{Modularity}: Each power system component (generator,
    demand, line, battery, reservoir, etc.) has its own LP assembly class
    (\texttt{*\_lp.cpp}) that independently adds variables, constraints,
    and objective terms to the shared \texttt{LinearProblem}.
  \item \textbf{Extensibility}: New components and solver backends can be
    added without modifying the core LP engine.  Solver backends are
    loaded as dynamic plugins (\texttt{libgtopt\_solver\_*.so}).
  \item \textbf{Interoperability}: JSON configuration with native Apache
    Arrow/Parquet I/O enables seamless integration with Python, R, and
    other data science ecosystems.
\end{enumerate}

\subsection{System Architecture}

Fig.~\ref{fig:arch_layers} illustrates the layered architecture.

\begin{figure}[t]
  \centering
  \begin{tikzpicture}[
    >=Stealth,
    layer/.style={draw, thick, rounded corners=3pt, minimum width=7cm,
      minimum height=1.1cm, font=\small, align=center},
    sublbl/.style={font=\tiny, text=black!60},
    arr/.style={->, thick, gray!70},
  ]
    % Layer 4 (bottom): Data Layer
    \node[layer, fill=green!12] (data) at (0,0)
      {\textbf{Data Layer}};
    \node[sublbl, below=-0.15cm of data]
      {JSON \quad Parquet/CSV \quad Schedule Resolution};

    % Layer 3: Model Layer
    \node[layer, fill=blue!10, above=0.5cm of data] (model)
      {\textbf{Model Layer}};
    \node[sublbl, below=-0.15cm of model]
      {Planning $\to$ System $\to$ Components};

    % Layer 2: LP Assembly Layer
    \node[layer, fill=orange!12, above=0.5cm of model] (lp)
      {\textbf{LP Assembly Layer}};
    \node[sublbl, below=-0.15cm of lp]
      {PlanningLP $\to$ SimulationLP $\to$ Component LPs};

    % Layer 1 (top): Solver Layer
    \node[layer, fill=red!10, above=0.5cm of lp] (solver)
      {\textbf{Solver Layer}};
    \node[sublbl, below=-0.15cm of solver]
      {CLP/CBC \quad CPLEX \quad HiGHS \quad (pluggable)};

    % Arrows between layers
    \draw[arr] (data.north) -- (model.south);
    \draw[arr] (model.north) -- (lp.south);
    \draw[arr] (lp.north) -- (solver.south);

    % Side labels
    \node[font=\tiny, rotate=90, anchor=south, gray]
      at (-3.9,1.2) {Data flow};
    \draw[->, thin, gray] (-3.9,0) -- (-3.9,3.6);
  \end{tikzpicture}
  \caption{\gtopt{} layered architecture: data ingestion, model
    construction, LP assembly, and pluggable solver backends.}
  \label{fig:arch_layers}
\end{figure}

The software is organized into four layers:

\begin{enumerate}
  \item \textbf{Data Layer}: JSON parsing (via DAW JSON Link, a
    zero-allocation header-only library), Parquet/CSV I/O (via Apache
    Arrow C++ SDK), and schedule resolution (\texttt{FieldSched<T>}
    variant type supporting scalar, vector, or filename references).
  \item \textbf{Model Layer}: The \texttt{Planning} $\rightarrow$
    \texttt{System} $\rightarrow$ component hierarchy that maps JSON
    input to C++ data structures.
  \item \textbf{LP Assembly Layer}: \texttt{PlanningLP} orchestrates
    \texttt{SimulationLP} instances (one per scenario/stage), which in
    turn assemble component LP objects (\texttt{GeneratorLP},
    \texttt{DemandLP}, \texttt{LineLP}, \texttt{BatteryLP},
    \texttt{ReservoirLP}, etc.) into a shared
    \texttt{FlatLinearProblem}.
  \item \textbf{Solver Layer}: The \texttt{SolverBackend} interface
    abstracts LP/MIP solution.  Implementations are loaded as dynamic
    plugins from configurable search paths.
\end{enumerate}

\subsection{Data Model and JSON Interface}

The input is specified as one or more JSON files containing three sections:

\begin{lstlisting}[language={},caption={Top-level JSON structure}]
{
  "options": { ... },
  "simulation": {
    "block_array": [...],
    "stage_array": [...],
    "scenario_array": [...]
  },
  "system": {
    "bus_array": [...],
    "generator_array": [...],
    "demand_array": [...],
    "line_array": [...],
    "battery_array": [...],
    ...
  }
}
\end{lstlisting}

Time-series data (demand profiles, generator capacity factors, reservoir
inflows) can be specified inline as JSON arrays, or as references to
external Parquet or CSV files in the \texttt{input\_directory}.  This
hybrid approach keeps the JSON configuration compact while supporting
arbitrarily large time-series via columnar Parquet storage.

\subsection{Python Toolchain}

\gtopt{} includes a comprehensive Python toolchain (16 registered CLI
tools):

\begin{itemize}
  \item \textbf{igtopt}: Excel workbook ($\rightarrow$ JSON + Parquet)
    converter for spreadsheet-based case preparation.
  \item \textbf{plp2gtopt}: PLP/PROMAX ($\rightarrow$ gtopt) converter
    for legacy hydrothermal scheduling data.
  \item \textbf{pp2gtopt / gtopt2pp}: pandapower
    ($\leftrightarrow$ gtopt) bidirectional converters for power flow
    model exchange.
  \item \textbf{gtopt\_compare}: Automated cross-validation against
    pandapower DC OPF on IEEE test networks.
  \item \textbf{gtopt\_diagram}: Network topology diagram generator
    using Mermaid syntax.
  \item \textbf{run\_gtopt / sddp\_monitor}: Execution management and
    SDDP convergence monitoring.
  \item \textbf{Validation tools}: \texttt{gtopt\_check\_json},
    \texttt{gtopt\_check\_lp}, \texttt{gtopt\_check\_output},
    \texttt{gtopt\_check\_solvers} for input validation and diagnostics.
\end{itemize}

\subsection{Web and GUI Services}

Two companion services provide graphical interfaces:

\begin{itemize}
  \item \textbf{webservice}: A Next.js REST API and browser UI for
    submitting optimization jobs, monitoring progress, and browsing
    results.
  \item \textbf{guiservice}: A Python/Flask application for creating
    and editing cases through a web-based form interface.
\end{itemize}

% ============================================================================
% Section VI: Validation and Benchmark Results
% ============================================================================

\section{Validation and Benchmark Results}
\label{sec:validation}

\gtopt{} is validated on a comprehensive suite of IEEE benchmark networks
and derived cases.  All test cases converge to optimal solutions (status~0)
with zero optimality gap.  Table~\ref{tab:benchmarks} summarizes the
benchmark suite.

\begin{table*}[t]
  \centering
  \caption{IEEE Benchmark Results (scale\_objective = 1000, CLP solver)}
  \label{tab:benchmarks}
  \begin{tabular}{@{}llcccccrc@{}}
    \toprule
    \textbf{Case} & \textbf{Type} & \textbf{Buses} & \textbf{Gen}
      & \textbf{Dem} & \textbf{Lines} & \textbf{Bat}
      & \textbf{Obj (scaled)} & \textbf{Blocks} \\
    \midrule
    \multicolumn{9}{@{}l}{\textit{Standard IEEE DC OPF (single snapshot)}} \\
    \texttt{ieee\_4b\_ori}   & DC OPF     & 4  & 2 & 2  & 5  & -- & 5.000    & 1  \\
    \texttt{ieee\_9b\_ori}   & DC OPF     & 9  & 3 & 3  & 9  & -- & 55.184   & 1  \\
    \texttt{ieee\_14b\_ori}  & DC OPF     & 14 & 5 & 11 & 20 & -- & 8.334    & 1  \\
    \texttt{ieee30b}          & DC OPF     & 30 & 6 & 21 & 41 & -- & 5.668    & 1  \\
    \texttt{ieee\_57b}       & DC OPF     & 57 & 7 & 42 & 80 & -- & 25.016   & 1  \\
    \midrule
    \multicolumn{9}{@{}l}{\textit{24-hour dispatch with solar profiles}} \\
    \texttt{ieee\_9b}        & Solar+OPF  & 9  & 3 & 3  & 9  & -- & 5400.007 & 24 \\
    \texttt{ieee\_14b}       & Solar+OPF  & 14 & 5 & 11 & 20 & -- & 1492.283 & 24 \\
    \midrule
    \multicolumn{9}{@{}l}{\textit{Battery energy storage}} \\
    \texttt{bat\_4b}         & BESS+OPF   & 4  & 3 & 2  & 5  & 1  & 14.160   & 4  \\
    \texttt{bat\_4b\_24}     & BESS+OPF   & 4  & 3 & 2  & 5  & 1  & 44.862   & 24 \\
    \midrule
    \multicolumn{9}{@{}l}{\textit{Capacity expansion}} \\
    \texttt{exp\_gen\_4b}    & Gen. exp.  & 4  & 2 & 2  & 5  & -- & 5.029    & 4  \\
    \texttt{exp\_bat\_4b}    & Bat. exp.  & 4  & 3 & 2  & 5  & 1  & 12.437   & 4  \\
    \texttt{c0}              & Dem. exp.  & 1  & 1 & 1  & 0  & -- & 23.163   & 5 stages \\
    \bottomrule
  \end{tabular}
\end{table*}

\subsection{IEEE 4-Bus Network}
\label{sec:ieee4b}

% ---- TikZ diagram: IEEE 4-bus ----
The IEEE 4-bus test case consists of a fully meshed network with 4 buses
and 5 transmission lines.  Two generators (g1: 300~MW at \$20/MWh,
g2: 200~MW at \$35/MWh) supply two demands (d3: 150~MW, d4: 100~MW).

Under DC OPF with Kirchhoff constraints, the cheaper generator g1
dispatches 250~MW and g2 idles.  The scaled objective value is
$5.0 = 250 \times 20 / 1000$.  Power routes through parallel paths
(b1$\rightarrow$b3 and b1$\rightarrow$b2$\rightarrow$b3) according to
the relative reactances.

\subsection{IEEE 9-Bus Network}
\label{sec:ieee9b}

The Anderson--Fouad 9-bus system~\cite{anderson2002power} has three
generation areas connected by a meshed transmission network.  Generators
g1 (250~MW, \$20/MWh), g2 (300~MW, \$35/MWh), and g3 (270~MW, \$30/MWh)
serve demands at buses 5, 7, and~9 (total 315~MW).

In the single-snapshot case (\texttt{ieee\_9b\_ori}), economic dispatch
follows merit order constrained by network topology.  The 24-hour variant
(\texttt{ieee\_9b}) replaces g3 with a solar plant (\$0/MWh) whose output
follows a daily capacity profile peaking at noon.  During daylight hours,
zero-cost solar generation displaces g2; at night, g2 carries the load.

\subsection{IEEE 14-Bus Network}
\label{sec:ieee14b}

The 14-bus system has 5 generators, 11 demands, and 20 transmission lines.
The constrained variant (\texttt{ieee\_14b}) tightens limits on lines
$\ell_{1-2}$ and $\ell_{1-5}$, forcing both to saturate simultaneously
at the evening demand peak.  This validates that shadow prices (Lagrange
multipliers of the binding Kirchhoff constraints) are computed correctly.

\subsection{IEEE 30-Bus and 57-Bus Networks}

The 30-bus~\cite{ieee30bus} and 57-bus~\cite{ieee57bus} systems test
scalability.  The 30-bus case has 6 generators, 21 demands, and 41 lines;
the 57-bus case has 7 generators, 42 demands, and 80 lines.  Both
converge to optimal solutions in under one second with the CLP solver.

\subsection{Battery Storage Cases}
\label{sec:bat_cases}

The \texttt{bat\_4b} case augments the 4-bus network with a 200~MWh
battery at bus~3 (60~MW charge/discharge, 95\% round-trip efficiency) and
a solar generator with a 4-period capacity profile.  The battery charges
during solar peak and discharges during evening demand, reducing total
system cost compared to the no-storage baseline.

The 24-hour variant (\texttt{bat\_4b\_24}) demonstrates full daily
charge/discharge cycling with state-of-charge tracking across all
24~blocks.

\subsection{Capacity Expansion Cases}
\label{sec:exp_cases}

Three expansion cases validate the investment planning formulation:

\begin{itemize}
  \item \texttt{exp\_gen\_4b}: Generator g1 can be expanded by up to
    290~MW at \$1000/year annualized cost.  The optimizer adds capacity
    to avoid dispatching expensive g2.
  \item \texttt{exp\_bat\_4b}: Battery expansion with 200~MWh modules
    at \$1000/year.  The optimizer installs storage to arbitrage
    solar/thermal cost differentials.
  \item \texttt{c0}: Multi-stage demand expansion (5 stages, 10\%
    discount rate).  Demand grows through 20~MW modules at
    \$8760/year.  The optimizer balances discounted investment cost
    against increasing served demand.
\end{itemize}

\subsection{Cross-Validation with pandapower}

The \texttt{gtopt\_compare} tool performs automated cross-validation
against pandapower~\cite{thurner2018pandapower} DC OPF on all IEEE test
networks (4, 9, 14, 30, 57 buses).  For each case, both tools solve the
same network data and compare:

\begin{itemize}
  \item Objective function value (total generation cost)
  \item Generator dispatch (MW per generator)
  \item Bus angles (radians)
  \item Line flows (MW)
  \item Locational marginal prices (bus balance duals)
\end{itemize}

Results match to within solver tolerance ($<10^{-6}$ relative error)
for all standard IEEE cases, confirming the correctness of \gtopt{}'s
DC OPF implementation.

% % ============================================================================
% Section VII: Case Studies (placeholder for user-defined cases)
% ============================================================================

\section{Case Studies}
\label{sec:cases}

% This section will be populated with user-defined case studies in a
% future revision.  Planned case studies include:
%
% \subsection{Chilean Hydrothermal System}
% Multi-year hydrothermal scheduling with SDDP, demonstrating hydro
% cascade modeling with reservoirs, turbines, waterways, and filtration.
%
% \subsection{Renewable Integration in a Distribution Network}
% Capacity expansion planning with high solar penetration, battery
% storage, and network congestion management.
%
% \subsection{Large-Scale Stochastic Expansion}
% Multi-stage, multi-scenario expansion planning demonstrating the
% cascade decomposition method on a 57-bus network with 10+ stages
% and 100+ scenarios.

% Uncomment when case studies are ready:
% \input{sections/case_study_hydro}
% \input{sections/case_study_renewable}
% \input{sections/case_study_expansion}

\textit{[Case studies with real-world applications will be included in a
future revision of this paper.  The benchmark results in
Section~\ref{sec:validation} demonstrate correctness and scalability on
standard IEEE test networks.]}
       % Placeholder for user-defined cases
% ============================================================================
% Section VIII: Conclusions and Future Work
% ============================================================================

\section{Conclusions}
\label{sec:conclusions}

A novel, high-performance tool for Generation and Transmission Expansion
Planning, designated \gtopt{}, has been presented.  The main conclusions
are as follows:

\begin{enumerate}
  \item \textbf{Complete formulation}: The GTEP problem is formulated
    as a sparse LP/MIP covering DC optimal power flow with Kirchhoff
    constraints, battery energy storage with state-of-charge tracking
    and self-discharge losses, hydro cascade systems with reservoirs,
    turbines, filtration, and volume-dependent efficiency, a generic
    water rights layer for irrigation agreement modeling with
    volume-dependent dynamic bounds, spinning reserve requirements,
    and multi-stage capacity expansion with annualized investment
    costs.

  \item \textbf{Three solution methods}: The monolithic LP provides
    exact solutions for moderate-size problems.  SDDP with
    multi-scenario apertures handles large stochastic problems
    efficiently, converging to within 1\% of the monolithic optimum
    in 10--15 iterations.  The novel cascade decomposition combines
    both approaches for long-horizon planning.

  \item \textbf{High performance}: The C++26 implementation with
    flat-map-based sparse-matrix assembly achieves a 10--27$\times$
    speedup over Python-based alternatives in LP assembly.  For the
    largest configuration tested (57-bus, 8760 blocks, ${\sim}18$M
    variables), assembly completes in approximately 8 seconds.
    Pluggable solver backends (CLP, CBC, CPLEX, HiGHS) are supported
    through a dynamic plugin architecture.

  \item \textbf{Validated correctness}: Cross-validation against
    pandapower DC OPF on IEEE benchmark networks (4--57 buses)
    demonstrates agreement to solver tolerance ($<10^{-6}$ relative
    error).  All 13 benchmark cases converge to optimal solutions.

  \item \textbf{Comprehensive toolchain}: A Python companion toolchain
    (igtopt, plp2gtopt, pp2gtopt, gtopt\_compare) and web interfaces
    provide a complete workflow from data preparation through result
    analysis, supporting interoperability with existing planning
    ecosystems.
\end{enumerate}

\subsection{Future Work}

Planned enhancements include:

\begin{itemize}
  \item \textbf{AC OPF support}: Extension of the formulation to include
    reactive power and voltage magnitude constraints via second-order
    cone programming (SOCP) relaxations.
  \item \textbf{Unit commitment}: Addition of binary start-up/shut-down
    decisions and minimum up/down time constraints for thermal generators.
  \item \textbf{Multi-sector coupling}: Extension of the model to include
    gas networks, district heating, and hydrogen production/storage for
    integrated energy system planning.
  \item \textbf{GPU-accelerated LP assembly}: Exploitation of GPU
    parallelism for sparse-matrix construction in very large-scale
    problems.
  \item \textbf{Machine learning integration}: Application of trained
    models to warm-start the SDDP algorithm or provide scenario
    reduction for stochastic programming.
  \item \textbf{Real-world case studies}: Application of \gtopt{} to
    national-scale expansion planning problems with detailed hydro
    cascade modeling.
\end{itemize}

The \gtopt{} tool is designed to support reproducible research in power
system optimization and to facilitate community collaboration in future
releases.


% ========== Acknowledgments ==========
\section*{Acknowledgments}
The authors gratefully acknowledge support from the Centro Nacional de
Inteligencia Artificial (CENIA, FB210017), funded by the Chilean National
Agency for Research and Development (ANID).
% Add further acknowledgments as needed.

% ========== References ==========
\bibliographystyle{IEEEtran}
\bibliography{references}

% ========== Biographies (optional for journal) ==========
% \begin{IEEEbiography}[{\includegraphics[width=1in,clip]{figures/matus.jpg}}]%
%   {Marcelo Matus~A.} received the ...
% \end{IEEEbiography}

\end{document}
